\documentclass[11pt]{article}
\usepackage{fontspec,amsmath,amssymb,amsthm,hyperref}
\usepackage[margin=1in]{geometry}
\setlength{\emergencystretch}{3em}
\title{Concentration around a minimizer without differentiability}
\author{LangevinLean mathematical exposition}
\date{}
\begin{document}\maketitle
\noindent Source module: \texttt{UnadjustedLangevin/GibbsConcentration.lean}. Compile this proof document with XeLaTeX or LuaLaTeX. The accompanying HTML pairs the mathematical argument with the exact Lean statements and records the scope of the formalization.


For a continuous \(q\)-strongly convex potential \(F\), \(q>0\), and an explicitly supplied minimizer \(z\), this module proves the equilibrium radius \(\int\|x-z\|^2d\pi_F\le2d/q\). The proof uses dilation and convexity, so it also applies to the nonsmooth potentials used by proximal algorithms.


\section{strongConvex mode quadratic bound}


Use strong convexity at the midpoint of \(z\) and \(x\), then use minimality \(F(z)\le F((z+x)/2)\). Rearranging yields \(F(x)\ge F(z)+(q/4)\|x-z\|^2\). This argument uses values of \(F\), not a derivative or a first-order optimality assertion at the minimizer.



\noindent Lean declarations: \texttt{strongConvex\_mode\_quadratic\_bound}.

\section{integrable exp neg of continuous strongConvex mode}


For \(q>0\), the midpoint bound gives \(e^{-F(x)}\le e^{-F(z)}e^{-(q/4)\|x-z\|^2}\). Translation invariance reduces integrability of the majorant to the centered Gaussian integral. Continuity makes the density measurable, so domination proves integrability without differentiability of \(F\).



\noindent Lean declarations: \texttt{integrable\_exp\_neg\_of\_continuous\_strongConvex\_mode}.

\section{integrable sq dist mul exp neg of continuous strongConvex mode}


Multiply the translated Gaussian bound by \(\|x-z\|^2\). Translation preserves Lebesgue measure, and a quadratic times a decaying Gaussian is integrable. Lean transports that integrability through \(x\mapsto x-z\) and then applies the nonnegative domination theorem to the desired weighted moment.



\noindent Lean declarations: \texttt{integrable\_sq\_dist\_mul\_exp\_neg\_of\_continuous\_strongConvex\_mode}.

\section{Contract toward the mode}


For \(0<r<1\), compare \(F(z+r(x-z))\) with \(F(x)\) using strong convexity and \(F(z)\le F(x)\). Exponentiating and using \(e^u\ge1+u\) produces a linear correction in \(\|x-z\|^2\). The dilation \(x\mapsto z+r(x-z)\) multiplies volume by \(r^d\), hence \(Z+(q/2)r(1-r)M\le r^{-d}Z\), with \(M=\int\|x-z\|^2e^{-F(x)}dx\).



\noindent Lean declarations: \texttt{gibbs\_mode\_dilation\_inequality}.

\section{integral sq dist mul exp neg le of continuous strongConvex mode}


Subtract \(Z\) from the dilation inequality and divide by \(1-r>0\). As \(r\uparrow1\), \((r^{-d}-1)/(1-r)\to d\), while \((q/2)rM\to(q/2)M\). The order-closed limit theorem yields \(M\le(2d/q)Z\). Lean proves the one-sided derivative limit of the scalar power before passing the inequality to the limit.



\noindent Lean declarations: \texttt{integral\_sq\_dist\_mul\_exp\_neg\_le\_of\_continuous\_strongConvex\_mode}.

\section{integrable sq dist gibbs of continuous strongConvex mode}


The density-measure integrability criterion changes the desired \(\pi_F\)-integral into a constant multiple of \(\int\|x-z\|^2e^{-F(x)}dx\). Its finiteness follows from the continuous, strongly convex mode estimate. No score moment or smoothness assumption is introduced during normalization.



\noindent Lean declarations: \texttt{integrable\_sq\_dist\_gibbs\_of\_continuous\_strongConvex\_mode}.

\section{integral sq dist gibbs le of continuous strongConvex mode}


First prove \(Z>0\), \(Z<\infty\), and integrability of the weighted distance. Then write \(\int\|x-z\|^2d\pi_F=M/Z\) and apply \(M\le(2d/q)Z\). Cancellation is justified by positivity of \(Z\); the result is \(\mathbb E_{\pi_F}\|X-z\|^2\le2d/q\), including \(d=0\).



\noindent Lean declarations: \texttt{integral\_sq\_dist\_gibbs\_le\_of\_continuous\_strongConvex\_mode}.

\end{document}
